\begin{theorem}[Tangent cohomology, direct images and miniversality]\hypertarget{tag-032U}{}\label{push:src:E.2.1}
\hypertarget{tag-01RH}{}\label{fam:src:A.5.2}\label{fam:st:A.5.2}\label{inv:D1.3}\label{inv:D2.1}\label{inv:D3.1}\label{cdp:st:E9}
\hypertarget{tag-02YO}{}\label{night:P5.3}\label{stmt:P5.3}
\hypertarget{tag-02YP}{}\label{night:P5.4}\label{stmt:P5.4}
\StatementRecord{push:src:E.2.1}{proved}{032U}{\cref{push:section:E.1,push:src:E.3.1,inv:D1.1,inv:D1.2}}
Let $k,l$ be positive integers with $\gcd(3k,4l)=1$, and let $X$
be a sphere member with multiplicities $(3k,4l)$ and the admissible
logarithmic data of \cref{push:section:E.1}. Put $z=kl$ and
$d=3k+4l$. On $\mathfrak B_{k,l}=\mathbb P(3k,4l)$, the vertical
root bundle $E$, the root lift $\bar f$ and its direct images satisfy
\[
 \begin{split}
 E&=\mathcal O\oplus\mathcal O(-z),\\
 R^q\bar f_*\mathcal O_X&=
 \mathcal O,\quad\mathcal O\oplus\mathcal O(-z),\quad\mathcal O(-z),\\
 K_X&=\bar f^*\mathcal O(z-d).
 \end{split}
\]
The tangent cohomology is
$\bigl(h^0,h^1,h^2,h^3\bigr)(X,T_X)=(1,1,1,0)$.
For the transition-arrow convention
$\kappa(\Phi)=(\partial_c\Phi)\Phi^{-1}$, the Kodaira--Spencer
class of $\partial_c$ is $e\otimes[\gamma]$, and the class of
$\partial_u$ is $-(2\pi i u)^{-1}e\otimes[\gamma]$.
Both classes are nonzero. The primary Kuranishi bracket vanishes,
the Kuranishi germ is a smooth reduced disc, and the fixed-data
period family is locally miniversal.

For the standard $(3,4)$ data, $e=\xi$ and the coarse direct
images are
\begin{equation}
 f_*T_X=\mathcal O\oplus\mathcal O(-1),\qquad
 R^1f_*T_X=\mathcal O\oplus\mathcal O(-1)^2,\qquad
 R^2f_*T_X=\mathcal O.
\label{pub:I:eq:D-4}
\end{equation}
\end{theorem}

\begin{proof}
\begin{enumerate}[label=\arabic*.,leftmargin=*]
\item The root lift, its proper flat Gorenstein charts, base change and
relative duality are those of \cref{push:src:E.3.1}.
Write $m=3k$ and $n=4l$. The finite root charts have the specified
free affine cyclic actions, and the cusp charts are $q=z_0z_1z_2$.

\item Use weighted coordinates $x,y$ of weights $m,n$, with coarse
coordinate $t=y^m/(y^m-x^n)$.
The two root divisors have bundles $\mathcal O(n)$ and
$\mathcal O(m)$, and the cusp has bundle $\mathcal O(mn)$.
The root-forgetting morphism
\[
 r:\mathfrak B_{k,l}\longrightarrow\mathfrak B_{1,1},
 \qquad (x,y)\longmapsto(x^l,y^k),
\]
uses the group homomorphism $\lambda\mapsto\lambda^{kl}$ on
the weighted quotient.  It is a holomorphic morphism of analytic
stacks over $B$, with local maps $s\mapsto s^k$ and
$s\mapsto s^l$.  In particular $r^*\mathcal O(1)=\mathcal O(z)$.
The analytic Picard group of this root curve has
the presentation
$\langle H,R_1,R_2\mid mR_1=H,\ nR_2=H\rangle$:
after subtracting the two stabilizer characters, a line bundle
descends to the coarse projective line.  Its only relations are
the two defining root identities.  The Smith form and
$\gcd(m,n)=1$ identify this group with $\mathbb Z$, with
$H=mn\mathcal O(1)$, $R_1=n\mathcal O(1)$, and
$R_2=m\mathcal O(1)$.  This proves the analytic torsion-free
Picard assertion used below.
Translations have identity vertical derivative and act trivially
on structure-sheaf cohomology.  Thus the vertical bundle computed
from the linear period factors is $r^*(\mathcal O\oplus
\mathcal O(-1))$.  The fixed vector gives a subline $\mathcal O$; its extension
splits because $H^1(\mathcal O(z))=0$.

\item Under the marked root-forgetting map in the preceding step,
the constant lattice cocycle $\gamma$ is unchanged. Its finite-chart
representatives and its cusp edge values $(0,-1,-1)$ therefore
give the nowhere-zero cohomology line proved in \cref{inv:D1.2}.
The cusp resolution and cup-product pairing are those of
\cref{push:src:E.3.1}; pulling the linear period factors through $r$
gives the quotient line $\mathcal O(-z)$.
Relative duality gives
$R^2\bar f_*\mathcal O_X=\mathcal O(-z)$ and
$\det R^1\bar f_*\mathcal O_X=\mathcal O(-z)$.
The extension of that quotient by the specified $[\gamma]$-line
splits because $H^1(\mathcal O(z))=0$, as computed in step~5.
Connected fibres give $\bar f_*\mathcal O_X=\mathcal O$.
These maps prove the three structure-sheaf direct images.

\item The logarithmic tangent sequence is
\[
 0\longrightarrow\mathcal V=\bar f^*E
 \longrightarrow\mathcal T=T_X(-\log W)
 \longrightarrow\mathcal H=\bar f^*\mathcal O(d-12z)
 \longrightarrow0.
\]
Its last map is the relative differential on smooth root charts
and sends $\sum a_i z_i\partial_{z_i}$ to
$(\sum a_i)q\partial_q$ at the cusp.  It is surjective there.
Since $K_{\mathfrak B}=\mathcal O(-d)$, determinants give
$K_X=\bar f^*\mathcal O(z-d)$.

\item The two quotient charts $x\ne0$, $y\ne0$ are finite quotients
of affine lines, and their overlap is a finite quotient of a
punctured line.  Finite averaging and Stein acyclicity compute
coherent cohomology by this two-chart complex.  Its homogeneous
Laurent monomials give
\[
 \begin{split}
 h^0\mathcal O(a)&=\#\{(i,j)\in\mathbb Z_{\ge0}^2:mi+nj=a\},\\
 h^1\mathcal O(a)&=h^0\mathcal O(-a-d).
 \end{split}
\]
The second equality identifies the remaining monomials
$x^{-i-1}y^{-j-1}$.
The bundles $\mathcal O(-z),\mathcal O(-2z)$ are acyclic.
Indeed a possible degree-one monomial would give
$m(i+1)+n(j+1)=az$, with $a=1$ or $2$.
Coprimality forces $i+1=lA$, $j+1=kB$, hence
$3A+4B=a$ with $A,B\ge1$, which is impossible.
Positive degree has no degree-one monomials, so
$H^1\mathcal O(z)=0$ as used above.

\item Put $b=d-12z$.  Both $b$ and $b-z$ are negative.
Degree-one monomials for $\mathcal O(b)$ would satisfy
$3A+4B=12$, with $A,B\ge1$, and none exist.
For $\mathcal O(b-z)$ the equation is $3A+4B=13$;
its only positive solution is $(A,B)=(3,1)$, producing
$x^{3l-2}y^{k-2}$ in the dual description.  It is allowed
exactly when $k>1$; write $e_k=1$ in that case and $e_k=0$
otherwise.  Projection formula and the two-column Leray sequence
now give
\[
 h^\bullet(\mathcal V)=(1,1,0,0),\qquad
 h^\bullet(\mathcal H)=(0,0,e_k,e_k),\qquad
 h^\bullet(\mathcal T)=(1,1,e_k,e_k).
\]

\item On the normalization $\nu:P=\mathrm{dP}_6\to W\subset X$,
the branchwise normal line is $K_P$: the invariant cusp form
$d\log x_1\wedge d\log x_2\wedge dq$ trivializes $K_X$
on the entire cusp neighbourhood.  Local normal coefficients
therefore give the exact sequence
\[
 0\longrightarrow\mathcal T\longrightarrow T_X
 \longrightarrow\nu_*K_P\longrightarrow0,
 \qquad h^\bullet(K_P)=(0,0,1).
\]
We must determine its connecting map when $k>1$.
The dual logarithmic sequence, tensored with $K_X$, has horizontal
subline $\bar f^*\mathcal O(13z-2d)$ and vertical quotient
$\bar f^*(\mathcal O(z-d)\oplus\mathcal O(2z-d))$.
The latter has no global sections, by the same equations
$3A+4B=1,2$ above.  Consequently
\[
 H^0(\Omega_X^1(\log W)\otimes K_X)
 =H^0(\mathfrak B,\mathcal O(13z-2d)).
\]
For $k>1$ this is generated by the monomial
$x^{3l-2}y^{k-2}$ multiplying the horizontal logarithmic
differential.  Both $x$ and $y$ are nonzero over $p_0$, so its
residue on $P$ is nonzero.  The ordinary logarithmic residue
sequence, tensored with $K_X$, identifies this residue map with
the Serre dual of
$H^2(P,K_P)\to H^3(X,\mathcal T)$.
One can check this identification directly in each chart
$q=z_0z_1z_2$: the residues of $dq/q$ are one on each branch,
and contraction with the invariant three-form is the normal
coefficient pairing.  Thus that connecting map has rank $e_k$.
The long exact sequence proves $h^2(T_X)=1$ and $h^3(T_X)=0$;
degrees zero and one are unchanged.

\item The derivative $\partial_c\Pi=e\otimes\gamma$ gives the constant
groupoid cocycle $\lambda\mapsto\gamma(\lambda)e$ on every
regular torus.  On the enlarged finite roots the derivative of the
affine generator is $(r_j/m_j)e$, and the derivative of its
punctured connector is $(\log s_j)/(2\pi i)\,r_je$.
The identity $m_j(r_j/m_j)=\gamma(w_j)$ verifies the power
relation, while $\gamma A_j=\gamma$ verifies conjugation with
every lattice arrow.  Changing $\log s_j$ by $2\pi i N$
adds the derivative of the integral arrow $Nw_j$; the connector
formula is used only on the punctured overlap.  Its central
root-chart representative is the constant $(r_j/m_j)e$,
which is holomorphic.  At the cusp $\gamma$ kills the toric periods,
so the constant deck cocycle is exactly its extending cocycle
above.  Its image under the two isomorphisms
$H^1(\mathcal V)\simeq H^1(\mathcal T)\simeq H^1(T_X)$
is $e\otimes[\gamma]\ne0$.  The chain rule gives the asserted
$u$-class.  With the tangent-lift convention $[\bar\partial h]$, this
class changes sign, as the explicit lift comparison in
\cref{fam:src:A.4.2} shows. The parameter and the transition-arrow
convention in the statement specify the displayed sign.

\item Let $\alpha$ denote the nonzero tangent class just constructed.
The holomorphic family gives a convergent deformation with this tangent.
Differentiating its integrability equation gives
$\bar\partial\alpha_2=-[\alpha,\alpha]/2$. Thus the primary
bracket vanishes; bilinearity and $h^1(T_X)=1$ give vanishing on
every pair.

By \cite[\S VII.8, principal theorem, pp.~598--599;
\S VIII.1, Theorem~1 and its corollary, p.~600]{Douady1974},
there is a minimal versal analytic germ with tangent space $H^1(T_X)$. The identification of this tangent space
uses tangent-valued \v Cech cocycles of transition maps, with changes
of trivialization acting by coboundaries. Its local algebra is
$\mathbb C\{t\}/I$. The parameter family induces
$t\mapsto g(c)$ with $g'(c_0)\ne0$. By the analytic inverse
function theorem, substitution by $g$ is injective; it annihilates
$I$, so $I=0$. The classifying map is an isomorphism of germs.
Pullback of the versal family realizes every sufficiently small
deformation and proves local miniversality.

\item For the standard data we compute the coarse direct images
by the following Koszul maps. All root bundles in this step are on
$\mathcal B=\mathbb P(3,4)$.
\begin{enumerate}[label=(\alph*),leftmargin=*]
\item Put $L=\OO_{\mathcal B}(-5)$. The first connecting map of \eqref{fam:eq:ab-log-tangent} is
the holomorphic bundle map
$\kappa:L\to E\otimes G$. In the splittings fixed in \ref{inv:D1.2} the
coefficient in $\OO\xi\otimes\OO\gamma$ lies in $H^0(\OO(5))=0$.
The remaining coefficients lie in
$H^0(\OO(4))=\C y$, twice, and $H^0(\OO(3))=\C x$.
Therefore, as a map from $E^*$ to $G$ with its $L$ factor understood,
\begin{equation}
 \kappa=\begin{pmatrix}0&a y\\ b y&c_0x\end{pmatrix}.
\label{pub:I:eq:D-5}
\end{equation}
We use $c_0$ for this coefficient to distinguish it from the deformation
parameter. The logarithmic period derivative at the cusp has leading
term the invertible matrix $B_0/(2\pi i)$. Hence
$\det\kappa=-ab y^2$ is nonzero there, so $ab\ne0$.
At the order-three point $\tau-\rho$ has order one on the root chart:
$j=1728t$ and the elliptic ramification order of $j$ at $\rho$ is
three. The marked torus period derivative is consequently nonzero
there: the two constant identity columns in the normalized period
matrix exclude an infinitesimal linear-frame change that could
otherwise conceal a nonzero derivative. In \eqref{pub:I:eq:D-5}, $y=0$ and
$x\ne0$ at that point, so $c_0\ne0$.
For the following matrices only, write the frames of $E,G$ as
$(e_1,e_2)$ and $(g_1,g_2)$ and replace them by
\[
 e'_1=(a/c_0)e_1,\quad e'_2=e_2,\qquad
 g'_1=b g_1,\quad g'_2=c_0g_2.
\]
The coefficients of $e'_1\otimes g'_2$, $e'_2\otimes g'_1$ and
$e'_2\otimes g'_2$ are then $y,y,x$.
The geometric cocycle $\xi\otimes\gamma$ keeps its original
meaning; $e'_1\otimes g'_1=(ab/c_0)\xi\otimes\gamma$ in these
matrix frames.

\item The next connecting map is cup product with $\kappa$. In suitable
bases it is
\begin{equation}
 L\otimes G=\OO(-5)\oplus\OO(-6)
 \xrightarrow{\ \delta_1\ }
 E\otimes\det G=\OO(-1)\oplus\OO(-2),\qquad
 \delta_1=\begin{pmatrix}-y&0\\-x&y\end{pmatrix}.
\label{pub:I:eq:D-6}
\end{equation}
The sign follows by wedging first with $\gamma$ and then with the
second generator of $G$; reversing both sign conventions leaves the
cokernels unchanged. As its determinant is $-y^2$, it is injective as
a morphism of sheaves. The first connecting map is
$\delta_0=(0,y,y,x)^t$ and is a subbundle inclusion, since $x,y$ have
no common zero. Its cokernel is
\begin{equation}
 \OO\oplus\OO(-1)\oplus\OO(2).
\label{pub:I:eq:D-7}
\end{equation}
Use the ordered frames
$u=e'_1\otimes g'_2+e'_2\otimes g'_1$ and
$v=e'_1\otimes g'_2$ for the two $\OO(-1)$ summands.
The change matrix has determinant $-1$, the image is $yu+xe'_2g'_2$,
and $v$ gives the unused summand. The remaining quotient is the exact sequence
$0\to\OO(-5)\xrightarrow{(y,x)}\OO(-1)\oplus\OO(-2)
\xrightarrow{(x,-y)}\OO(2)\to0$.

\item The cokernel of \eqref{pub:I:eq:D-6} is supported at the order-three root point. There $x$ is a unit. For target generators $q_1,q_2$, the
relations are $-yq_1-xq_2=0$ and $yq_2=0$. Eliminating
$q_2=-(y/x)q_1$ gives the single relation $y^2q_1=0$ and hence
\begin{equation}
 Q=\OO(-1)\big|_{2p_1},
\label{pub:I:eq:D-8}
\end{equation}
where $2p_1$ is the length-two divisor in the root coordinate $y$. The generator is the first target component,
and its annihilator is $y^2$. The long direct-image sequence now gives
\[
 R^0\bar f_*\mathcal T=E,\quad
 R^1\bar f_*\mathcal T=\OO\oplus\OO(-1)\oplus\OO(2),
\]
\begin{equation}
 0\longrightarrow Q\longrightarrow R^2\bar f_*\mathcal T
 \longrightarrow\OO(-6)\longrightarrow0.
\label{pub:I:eq:D-9}
\end{equation}
This last extension splits, since $\operatorname{Ext}^1(\OO(-6),Q)
=H^1(Q(6))=0$.

\item Let $\nu:P=\mathrm dP_6\to W\hookrightarrow X$ be the normalization
immersion. There is an exact sequence
\begin{equation}
 \hypertarget{tag-01P1}{}\label{fam:eq:ab-normal-quotient}
0\longrightarrow\mathcal T\longrightarrow T_X
 \longrightarrow\nu_*K_P\longrightarrow0.             
\end{equation}
To verify the quotient, in a chart $W=(z_1\cdots z_r=0)$ the map takes
a vector field to its normal coefficient on each of the $r$ separate
branches. The coefficients lift independently; its kernel is precisely
the vector fields tangent to every branch. Thus the normal-coefficient map identifies the quotient with
the pushforward of the normal line of $\nu$.
The form $d\log x_1\wedge d\log x_2\wedge dq$ becomes a unit times
$dz_1\wedge dz_2\wedge dz_3$ in every unimodular chart and is invariant
under every cusp deck transformation. It trivializes $\nu^*K_X$.
Adjunction for the immersion gives its normal line $K_P$, proving
\eqref{fam:eq:ab-normal-quotient}. Since $P$ is rational, its $K_P$ cohomology is $(0,0,1)$ by
Serre duality. Finite pushforward preserves that cohomology.

\item It remains to identify the extension at $p_0$ in degree two. We give
a fibre-sensitive calculation. The kernel of
\[
 \Omega_X^1|_W\longrightarrow\nu_*\Omega_P^1
\]
is $\nu_*\OO_P\otimes N_W^*$. At $xy=0$ its sections are
$a(y,z)y\,dx+b(x,z)x\,dy$; at $xyz=0$ they are
$a(y,z)yz\,dx+b(x,z)xz\,dy+c(x,y)xy\,dz$.
These identifications glue by the common conormal $dq$.
Since $N_W$ and $K_X|_W$ are trivial and $H^0(P,\Omega_P^1)=0$, the
Serre-dual space of $H^2(W,T_X|_W)$ is one-dimensional.
Top-degree base change identifies this group with the residue fibre
of $(R^2\bar f_*T_X)_{p_0}$. Its generic rank is one by \eqref{pub:I:eq:D-9}.
A finite module over the discrete valuation ring $\C\{q\}$ of rank
one and residue dimension one is free of rank one. Thus the extension
by $\C_{p_0}$ increases the degree of its line part by $p_0$, without
adding torsion. We obtain the complete root direct images
\[
 R^0\bar f_*T_X=\OO\oplus\OO(-1),\quad
 R^1\bar f_*T_X=\OO\oplus\OO(-1)\oplus\OO(2),
\]
\begin{equation}
 R^2\bar f_*T_X=\OO(6)\oplus Q.
\label{pub:I:eq:D-11}
\end{equation}
On the order-three chart, an invariant coefficient of the generator
of $Q$ must have exponent $2$ modulo $3$, whereas $Q$ only retains
exponents $0,1$. Therefore $\rho_*Q=0$. Applying the exact $\rho_*$
to \eqref{pub:I:eq:D-11} and using
$d(0)=0,d(-1)=d(2)=-1,d(6)=0$ proves \eqref{pub:I:eq:D-4}.
The coarse direct images are the ones displayed in the statement;
their Leray cohomology agrees with the tuple obtained in step~7.
\end{enumerate}
\end{enumerate}
\end{proof}
